\begin{proposition}\label{prop:decic-no63}
An elliptic curve over a number field of degree $10$ has no point
of order $63$.
\end{proposition}
\begin{proof}
Put $Y=X_1(63)/\langle8\rangle$, of genus $49$.
The involution acts freely on the source also in characteristic
$2$, including at the cusps, so $Y$ has good reduction at $2$.
The complete invariant-character winding calculation accounts for
all $49$ dimensions and proves $J_Y(\mathbb Q)$ has rank zero
by the winding criterion of \cite[Section~3]{DEHMZ}.
We use the direct model obtained from
\[
 c=t^2(t-1),\quad b=c(t^2-t+1),\qquad
 E_9:y^2+(1-c)xy-by=x^3-bx^2
\]
by adjoining $u=x(P_7)$ with $P_7$ of order $7$.
The polynomial $\psi_7(E_9)(u)$ is irreducible of degree $24$;
the resulting function field has genus $49$ and $j$-degree $864$.
Here $(0,0)$ is the marked point of order $9$, and the quotient
forgets the sign of the $7$-point.

Let $D_X$ be the trace divisor of a hypothetical level structure
over a degree-$10$ field, and let $D$ be its pushforward to $Y$.
The complete source noncuspidal place counts at $2$, in degrees
$1,\ldots,10$, are
\[
                    (0,0,0,0,0,0,54,54,72,54).
\]
Hence a noncuspidal reduction has one source place of degree
$7,8,9$ or $10$ and a cuspidal complement.
Their images on $Y$ have the same degrees and form respectively
$3,3,4,3$ diamond orbits. This is checked by exhaustive enumeration
of the marked pairs $(t,u)$ and their Frobenius orbits.

Use point-level diamonds $\delta_a$ and put
$A_5=T_5-1-5\delta_5$. Eichler--Shimura and rational torsion
injectivity at $5$ give $A_5J_Y(\mathbb Q)=0$.
A rational unit-component cusp $c_0$ satisfies $A_5c_0=0$.
The finite cusp signature is $1^9 2^3 3^9 6$, and the global
signature is $1^9 2^3 3^6 6\,9$.
Thus every complement of degree at most $2$ lifts globally.
For such a complement, applying $A_5$ to $[D-10c_0]$ makes
$A_5P$ principal, where $P$ is the noncuspidal residual place.
The accepted exact tests exclude the degree-$8$ and degree-$9$
orbits, and the fresh tests give $h^0(A_5P)=0$ on all three
new degree-$10$ orbits.

For a degree-$7$ place the complement has degree $3$.
Every complement except one of the three cubic places obtained
by splitting the global degree-$9$ cusp lifts globally, and the
accepted degree-$7$ tests exclude it.
The three remaining cubic places, say $C_1,C_2,C_3$, are exactly
the cubic poles of $j$ of width $1$.
On this boundary, cyclotomic Galois $G_5$ equals the point diamond
$\delta_5$, and the Hecke formula gives
\[
 T_5C_i=G_5C_i+5C_i,\qquad
                   A_5C_i=4(C_i-G_5C_i).
\]
Primitive Tate-column arithmetic shows that $G_5$ cycles the three
closed cubic places. We test all six ordered pairs $i\ne j$,
so no orientation choice enters the argument.
For each of the three degree-$7$ representatives the six exact
computations give
\[
                      h^0\bigl(A_5P+4(C_i-C_j)\bigr)=0.
\]
These $18$ tests exclude the remaining complements.
The actual six geometric $5$-kernels form one orbit of length $6$
in degree $7$, and two orbits of length $3$ in degree $10$.
Every norm multiplicity is $1$.
A fresh native reconstruction verifies all $21$ new operator
divisors and their zero Riemann--Roch spaces, using exact
two-generator prime ideals. Independent marked Velu maps check
all $36$ kernels, both marked points, and every residue norm.

It remains to exclude entirely cuspidal reduction.
Put $U=(\mathcal C-1)H^0(Y_{\mathbb F_2},\Omega^1)$,
where $\mathcal C$ is Cartier. The exact computation gives
$\dim U=45$. The retained restriction kernels through weighted
degree $9$ cover $239226$ degree-$10$ cuspidal divisors;
every restriction from $U$ has rank $10$.
The only missing configurations are $10c$ at the nine rational
cusps and $5c$ at the three quadratic cusps.
Fresh differential-space kernels give rank $10$ in all twelve
cases. Each uses its own full differential basis, so no bases
from different exports are combined.
A separate complete-matrix enumeration checks all $239238$
cuspidal configurations and the twelve new ranks.

The full-Jacobian dyadic criterion
(Corollary~\ref{cor:dyadic-Cartier-formal-immersion}) now applies.
The entire closed-cusp trace lifts over $\mathbb Q_2$, including
at a cubic cusp which does not lift globally.
Rank zero and Manin--Drinfeld make $[D-C]$ torsion over
$\mathbb Q_2$, and $T_2-\delta_2$ kills its reduction kernel.
Its cotangent image is $U$, whose restrictions are all surjective.
Formal immersion therefore gives $D=C$, contrary to noncuspidal
support.
\end{proof}
